% Transcription — fonds Alexandre Grothendieck, shelfmark 22, batch 2 (pages 21-29).
% Established from the facsimile archives/batches/22/batch-02.pdf (a mirror of
% https://grothendieck.umontpellier.fr/22.pdf), per the skill
% transcribe-grothendieck. The batch file carries no cover sheet and opens at
% page 21: PDF page N is archivists' page 20+N. Checked against the pencilled
% numbers bottom-left, which read 21, 22, 23, 24, 25, 26, 28 and 29 on the
% corresponding leaves; page 27 carries its number in the opposite corner because
% the leaf was filmed inverted. This batch is the tail of the folder — nine
% leaves, not twenty.
%
% Pass: Opus 5 (claude-opus-5), 2026-09-19 — first pass, unchecked against the pages by a human.
%
% The sorting, which the folder's three kinds of paper force page by page:
%
% - Pages 27 and 29 are printed matter — a reproduced typescript, in a type that
%   is not his and in a hand that is nobody's, on the topology of the space of
%   répartitions and the group of idèles, norms and traces along an extension of
%   function fields. It is unrelated to chapter 0_IV and carries no annotation of
%   his: two pencilled archivists' numbers and nothing else. Under the rule for
%   offprints it is neither transcribed nor described, and pages 27 and 29 simply
%   do not appear below. (They are the versos of the two pencil leaves 26 and 28,
%   which is why they fall where they do.)
%
% - No page of this batch is typescript. The "tapuscrit annoté" of the inventory
%   is not in pages 21-29; whatever it is, it lies in batch 1.
%
% - The seven remaining pages are manuscript, in two distinct campaigns:
%     21-24  ink, on the back of letterhead sheets of his typist (M. Dufour,
%            varitypiste). One continuous run, unpaginated by him.
%     25     an otherwise blank cream sheet carrying only « Différentielles » in
%            his hand, top right, underlined. It is a cover for the run that
%            follows, it takes no \page{}, and its word opens the section below.
%            Page 21 already refers forward to it: « cf. "Diff" ».
%     26, 28 pencil, a different and later-looking campaign, on cream paper.
%
% What the batch is about. Two questions, both from chapter 0_IV.
%
%   21-24  Is A = K[[t_1,…,t_n]] formally smooth over k, and for which topology?
%          Page 21 divides into cases on the characteristic and settles p > 0 by
%          p-bases (a p-base of K over k, completed by the t_i, is a p-base of A
%          over k). Page 22 turns to p = 0 and announces the opposite: « je suis
%          pessimiste », A = k[[t]] is never formally smooth over k for the
%          discrete topology. The obstruction is exhibited: with k = Q((x_i)),
%          Ω_k ⊗_k A ≃ A^(I) sits strictly inside ∏_i A dx_i, so a derivation
%          D : A → Ω_k ⊗_k A extending d_{K/k} would have to land in a direct sum
%          while D(Σ a_λ t^λ) = Σ d_{K/k}(a_λ) t^λ has infinite support.
%          Page 23 states this as a Proposition (Ω_{K/k} of finite dimension over
%          K) and page 24 runs the construction of the counterexample f = Σ a_n t^n.
%
%   26, 28 A « Remarque » on the square k → K, k ∩ K^p → K^p, and the question
%          whether Ω^1_{k/k∩K^p} ⊗_k K → Ω^1_K is injective — then, on page 26's
%          lower half and all of page 28, an explicit counterexample in
%          characteristic p built over a field F with x, y and a_1,…,a_{p-1}
%          algebraically independent, y = Σ_0^{p-1} a_i x^i, a = x^p.
%
% Legibility. Pages 21-24 are a fast draft, heavily struck and abbreviated: the
% formulas carry, the connective prose largely does not, and the apparatus below
% records that honestly rather than smoothing it. Pages 26 and 28 are pencil on
% cream paper with a typed verso showing through, and are harder still. No date
% appears anywhere in these nine leaves.
\documentclass[11pt,a4paper]{article}
% Le préambule commun à toutes les transcriptions du fonds Grothendieck.
%
% Il tient en une idée : **l'incertitude se compose, elle ne se résout pas.**
% Une transcription qui lisse un mot douteux a détruit la seule chose pour
% laquelle on la faisait — un lecteur doit pouvoir savoir, à chaque endroit,
% ce qui était sur le feuillet et ce qui a été ajouté. Les six macros qui
% suivent sont donc l'appareil critique, et non de la décoration.
%
% Les mêmes commandes sont reconnues par `scripts/render.mjs`, qui produit la
% vue de lecture du site. Une macro ajoutée ici doit l'être aussi là-bas,
% faute de quoi le rendu échouera bruyamment — ce qui est voulu : un rendu
% silencieusement faux est pire qu'un rendu absent.

\ProvidesPackage{grothendieck}[2026/08/08 Transcriptions du fonds Grothendieck]

\usepackage{amsmath,amssymb,amsthm}
\usepackage{mathrsfs}
\usepackage{stmaryrd}
% Les diagrammes commutatifs sont partout dans ce fonds : c'est la langue
% même de la géométrie algébrique de Grothendieck, et une transcription qui
% les rendrait en prose ne transcrirait rien.
\usepackage{tikz-cd}
% Français par défaut — c'est la langue des feuillets — mais l'anglais est
% chargé aussi : la traduction et le résumé sont des documents anglais, et la
% typographie française (espaces avant les deux-points) y serait une faute.
% Ces documents basculent par \selectlanguage{english} en tête de corps.
\usepackage[english,french]{babel}
\usepackage{fontspec}
\usepackage{geometry}
\geometry{a4paper,margin=25mm,marginparwidth=28mm}
\usepackage{marginnote}
\usepackage{xcolor}
% ulem, not soul. `soul`'s \st and \ul are built on a character-by-character
% scan that XeLaTeX's font machinery breaks: the document fails with an
% undefined control sequence rather than degrading. ulem does the same two jobs
% and is Unicode-engine safe. `normalem` stops it hijacking \emph, which the
% transcriptions use for Grothendieck's own emphasis.
\usepackage[normalem]{ulem}
\usepackage{hyperref}
\hypersetup{colorlinks=true,linkcolor=black,urlcolor=black,pdfborder={0 0 0}}

% --- Métadonnées du lot -----------------------------------------------------
% Reprises telles quelles par le script de rendu, qui les affiche en tête de
% la vue de lecture. Elles fixent ce que le fichier prétend couvrir : sans
% elles, une transcription flotte sans point d'attache dans un fonds de seize
% mille feuillets.

\newcommand{\folder}[1]{\gdef\grothFolder{#1}}
\newcommand{\batch}[1]{\gdef\grothBatch{#1}}
\newcommand{\leaves}[2]{\gdef\grothFirst{#1}\gdef\grothLast{#2}}
\newcommand{\foldertitle}[1]{\gdef\grothFoldertitle{#1}}
\gdef\grothFolder{?}\gdef\grothBatch{?}\gdef\grothFirst{?}\gdef\grothLast{?}\gdef\grothFoldertitle{}

% --- Le résumé --------------------------------------------------------------
%
% Il ouvre la lecture modernisée, sous le titre « Résumé », et il est composé
% en petit. Ce n'est pas un ornement : c'est le seul passage du document qui
% ne soit pas des mathématiques, et un lecteur doit pouvoir le reconnaître
% avant de l'avoir lu — puis le sauter s'il connaît déjà le sujet, ou s'y
% arrêter s'il ne le connaît pas. Le filet qui le suit marque la reprise du
% corps.

\newenvironment{resume}
  {\small\setlength{\parskip}{0.45em}}
  {\par\medskip\hrule\medskip}

% --- Mots-clés ---------------------------------------------------------------
%
% En anglais, en clôture du résumé de la lecture modernisée : le vocabulaire
% moderne sous lequel chercher ce que les feuillets construisent. Le manifeste
% du site (`npm run manifest`) les extrait pour étiqueter la cote entière —
% cette ligne est la source unique des tags, il n'y a pas de fichier à côté.
\newcommand{\keywords}[1]{\par\smallskip\noindent
  {\footnotesize\textsc{Keywords} --- \textit{#1}}\par}

% --- L'appareil critique ----------------------------------------------------

% Le numéro de feuillet, en marge. C'est l'ancre qui permet de régler un
% désaccord sur une formule en regardant *un* feuillet plutôt qu'une chemise
% de six cents. Le site s'en sert aussi pour tourner les pages du fac-similé
% quand on fait défiler la transcription.
\newcommand{\leaf}[1]{%
  \marginnote{\footnotesize\color{gray!70}\textsf{#1}}%
  \label{leaf:#1}\ignorespaces}

% Illisible : on ne devine pas. Un mot inventé qui se lit comme les autres est
% le pire résultat possible.
\newcommand{\ill}{\textcolor{red!60!black}{[\,\ldots\,]}}

% Lecture douteuse : le mot est donné, et signalé comme incertain.
\newcommand{\uncertain}[1]{\uline{#1}}

% Ajout de l'éditeur — une lettre manquante, un mot sous-entendu.
\newcommand{\add}[1]{\textcolor{blue!50!black}{[#1]}}

% Biffé par Grothendieck lui-même. Ce qu'il a rayé fait partie du document :
% une rature dit souvent où la pensée a changé de direction.
\newcommand{\struck}[1]{\sout{#1}}

% Note du transcripteur, sur l'état matériel du feuillet.
\newcommand{\note}[1]{\par\smallskip{\small\color{gray!60!black}\textit{[#1]}}\par\smallskip}

% Note marginale de Grothendieck, distincte du corps du texte.
\newcommand{\marginal}[1]{\par\smallskip\noindent\hspace{1em}%
  {\small\color{gray!60!black}#1}\par\smallskip}

% --- Titre ------------------------------------------------------------------

\AtBeginDocument{%
  \begin{center}
    {\large\bfseries Fonds Alexandre Grothendieck --- cote n\textsuperscript{o}~\grothFolder}\\[2pt]
    {\itshape\grothFoldertitle}\\[4pt]
    {\small feuillets \grothFirst--\grothLast\ (lot \grothBatch)}\\[2pt]
    {\footnotesize\color{gray!60!black}Transcription établie d'après le fac-similé numérisé
     par l'Université de Montpellier.\\
     Les lectures incertaines sont soulignées, les illisibles notées [\,\ldots\,],
     les ajouts entre crochets.}
  \end{center}
  \vspace{1em}}

\endinput


\folder{22}
\batch{2}
\pages{21}{29}
\dating{1963-[vers 1967]}
\watermark{Édition de démonstration}
\foldertitle{[EGA IV]. [Chapitre] 0 IV. Cohen, différentielles, algèbres formellement simples (Compléments) : tirés à part (1963), tapuscrit annoté (s.d.), notes manuscrites (s.d.).}

\begin{document}

\page{21}
\uncertain{loc. artin. complets}, $A$ séparé \ill{} \ill{} \ill{} pour la
topologie habituelle, alors $A$ est-il $=$ \ill{} \ill{} pour la topologie
discrète ???
\note{les premières lignes du feuillet sont une question qu'il se pose à
lui-même ; elles ne se laissent lire qu'en partie, et
« \uncertain{loc. artin. complets} » est écrit en interligne au-dessus de la
ligne}

1°) \emph{Cas $p = 0$}. On se ramène au cas où $k$ est le corps premier
\ill{} \ill{} \struck{\ill{}} \ill{} $K/k$ une extension, $K/k$ l'ex.
considérée.

2°) \struck{\emph{Cas $p \geq 0$}} Alors on a $A \simeq K[[t_1,\dots,t_n]]$,
$K/k$ une extension séparable. \struck{\ill{}} \ill{} $k$, il
\struck{\ill{}} \ill{} que $A$ \ill{} sur $K$.

Supposons que $K/k$ \uncertain{admet} une p-base
\struck{\ill{} $K$ est parfait, alors}

3°) \emph{Cas $p > 0$}. \struck{Alors $A$ est une \ill{}}
\ill{}, i.e.\ \ill{} $(t_1,\dots,t_n)$
\ill{} $A$, ayant une p-base sur $K$, \ill{} \struck{\ill{}}

\marginal{$\Omega_{K/k}$ \uncertain{est libre} sur $K$.}

p-base de $A$ sur $k$. Il me semble \uncertain{si}
$K$ est fini sur $K^{p}$, on prend une p-base $(x_1,\dots,x_n)$ de $K$, et on
y ajoute $t_1,\dots,t_n$. \uncertain{Vu} la complétion pour les $t_i$, on
trouve une p-base de $A$ sur $k$, et $k(A^{p} \simeq \uncertain{K}
\otimes_{K^{p}} A^{p}$.
\note{la parenthèse ouverte après $k$ n'est pas refermée sur le feuillet ; le
terme après $\simeq$ est repassé et ne se lit pas avec certitude}
[cf. « Diff »]. Donc dans ce cas \ill{} \ill{}
\note{le renvoi vise le feuillet de couverture « Différentielles » qui ouvre la
section suivante, quatre pages plus loin}

\page{22}
\struck{3°)} \emph{Cas contraire}. Je suis \uncertain{pessimiste}.

1°) \emph{Cas $p = 0$}. \ill{} \ill{} $k[[t]]$ sur $k$, \ill{} \ill{}
\ill{} \ill{} Je présume que de dim $A \geq 0$, $A$ n'est jamais
\uncertain{formellement lisse} sur $k$ pour la topologie discrète.

Donc, dans les cas défavorables, je \uncertain{cherche} \ill{} qu'il y a
\struck{\ill{}} \ill{} naturellement un corps \ill{} [Cf.\ \ill{} \ill{}
\ill{} \ill{} \ill{}].

\[
  \Omega^{1}_{k} \otimes_{k} A \longrightarrow \Omega^{1}_{A}
\]
\struck{\ill{} être un iso ?, pourquoi ?}

\emph{Supposons} \struck{\ill{} $\Omega_{k}$} \quad
$k = \mathbf{Q}((x_i))$ \quad \struck{$k_0 = \mathbf{Q}$}

$A = k[[t]]$, \quad
$\Omega_{k} \otimes_{k} A \simeq A^{(I)} \subsetneq \prod_{i} A\,dx_i$,
\note{le symbole de produit est écrit par-dessus une rature}
\ill{} cherchons
$A \xrightarrow{\;D\;} \Omega_{k} \otimes_{k} A$, qui soit donnée sur la base
\[
  D(x_i) = 1/\alpha\; x_i .
\]
\note{la lettre au dénominateur est une seule lettre, lue $\alpha$ sans
certitude ; elle revient deux lignes plus bas}
Si \uncertain{l'on a} une telle $D$, alors $D(a)$ \ill{} \ill{} combinaison
\uncertain{linéaire} finie de certains $1/\alpha\, x_i$, $i \in J$, \ill{} et
posant $\mathbf{Q}$ par $k_0 = \mathbf{Q}((x_i : i \in J))$, \ill{} \ill{}
pour \ill{} que $dt = 0$. \struck{\ill{}}

\page{23}
\emph{Proposition.} Soient $k$ un corps, $K$ une extension telle que
$\Omega_{K/k}$ \uncertain{soit} de dim finie sur $K$, et
$A = K[[t_1,\dots,t_n]]$ \struck{\ill{}} de dim finie sur $K$,
\ill{} $K$, \ill{} $n \geq 0$. Alors $A$, \ill{} \ill{} \ill{}
\uncertain{discret}, \ill{} \ill{} \ill{} sur $K$ rel.\ \ill{} ; $k$, et
\ill{} : donc \ill{} \ill{} \ill{}
\note{l'énoncé est écrit d'un trait, sans reprise, et sa moitié droite ne se
laisse pas lire ; ce qui est sûr est l'hypothèse — $\Omega_{K/k}$ de dimension
finie sur $K$ — et l'anneau $A = K[[t_1,\dots,t_n]]$}

Je \ill{} \uncertain{pour} \ill{} que l'\ill{}
\[
  \Omega_{K/k} \otimes_{K} A \longrightarrow \Omega_{A/k}
\]
\ill{} par \ill{} : \uncertain{quotients}, i.e.\ \ill{} \ill{}
\uncertain{distinct}. \ill{}
\[
  D : A \longrightarrow \Omega_{K/k} \otimes_{K} A
\]
\ill{} \uncertain{induisant} $d_{K/k} \otimes_{K} A$. \ill{}
\uncertain{une telle} $D$, \ill{} \ill{} \ill{} que les $D(t_i)$
\uncertain{soient} dans \ill{} \ill{} $V \otimes_{K} A$, avec
$V \subset \Omega_{K/k}$ \ill{} \ill{} \ill{} il existe \ill{} \ill{}
\uncertain{sous-corps} $K'$ de $k$, \ill{} \ill{} \ill{} de type fini
\ill{} $I \subset \Omega_{K'/K} \otimes_{K} A$, \ill{}
$\Omega_{K/k} \otimes_{K} A$ \ill{} \ill{} \ill{} \uncertain{contenue} dans
\emph{\uncertain{type fini}}, \ill{} \ill{} \ill{} $dt_i$ \ill{} \ill{}
\ill{} $A$, donc \ill{} \ill{} \ill{} \ill{} \uncertain{l'hypothèse !}
\note{tout le bas du feuillet est un palimpseste rapide ; seuls les objets
nommés — $V$, $K'$, les $D(t_i)$ — se laissent relever}

\page{24}
\ill{} $D t_i = 0$. \struck{\ill{}} Prenons \ill{} \emph{base} \ill{}
$\Omega_{K/k}$ \ill{} $K$, \ill{} $(\omega_i)_{i \in I}$, donc
$\Omega_{K/k} \otimes_{K} A = A^{(I)} \subset A^{I}$. \uncertain{En} \ill{}
\ill{} \ill{} \ill{} $A^{I}$, $D$ \uncertain{est donnée} \ill{} \ill{}
\ill{}, $D_i$, \ill{} \ill{} \ill{} \uncertain{uniquement déterminé} \ill{}
\[
  D\Big(\sum_{\lambda} a_{\lambda} t^{\lambda}\Big)
  \;=\; \sum_{\lambda} d_{K/k}(a_{\lambda})\, t^{\lambda} .
\]
\note{le $d$ du membre de droite est repassé plusieurs fois}

Il \uncertain{veut} \ill{} \ill{} qui \ill{} \ill{} \ill{}
$t = \sum a_{\lambda} t^{\lambda}$ \ill{} \ill{} \ill{}
$(D\varphi)_{\lambda \in I}$ \ill{} $\notin A^{(I)}$.
\struck{Je \ill{}} Choisissons \ill{} \struck{\ill{} $D(a_{\lambda})$
\ill{}} \ill{} \emph{\ill{}} $i_0, i_1, \dots, i_{\ill{}}$ \ill{} $i$,
\struck{telle que} \ill{} \ill{} d'où \ill{}
\struck{$D_{i_n}(a_n) \neq 0$} \ill{} \ill{} \ill{} \ill{} $a_n \in K$
\ill{} \ill{} que \struck{$d_{K/k}\, a_n \neq 0$)} il \ill{} \ill{}
\ill{} prendre
\[
  f = \sum a_n t^{n} . \qquad \dots
\]
\note{le feuillet s'arrête sur cette construction : la suite du contre-exemple
n'est pas sur la page}

\section{Différentielles}
\note{le mot est de sa main, souligné, seul en haut à droite du feuillet 25,
autrement blanc : c'est une couverture pour ce qui suit, et c'est à elle que
renvoie le « cf. "Diff" » du feuillet 21. Le feuillet ne porte rien d'autre et
ne reçoit pas de \texttt{page}}

\page{26}
\emph{Remarque} \quad Soient \ill{}

\begin{tikzcd}
k \arrow[r] \arrow[d, no head] & K \arrow[d, no head] \\
k' = k \cap K^{p} \arrow[r, no head] & K^{p}
\end{tikzcd}

\note{les deux traits verticaux et le trait du bas sont sans pointe sur le
feuillet ; seule la flèche du haut en porte une}

dans \ill{} il \ill{} \struck{\ill{}} $\gamma$ de $k$ \uncertain{provenant}
\add{dans $K$} être \uncertain{principalement} \struck{\ill{}}
$k \cap K^{p}$, et
\note{les crochets d'ajout marquent ici, et une fois encore au feuillet 28, un
ajout de sa main porté en interligne, non une addition de l'éditeur} \ill{} (\ill{} $K^{p}$, \ill{} \ill{} \ill{}
i.e.\ $d_{k/k'}\, x$ et $d_{K/k'}\, \gamma$ \uncertain{proviennent} de
\ill{} \uncertain{indiqués}, \ill{} \ill{} \ill{} \ill{}
$d_{K/K^{p}}\, x$ et $d_{K/K^{p}}\, \gamma$ \ill{} \ill{} ;
\ill{} \ill{} \ill{} \ill{} injection,
\[
  \Omega^{1}_{k/k \cap K^{p}} \otimes_{k} K \longrightarrow \Omega^{1}_{K}
\]
\ill{} \ill{} \ill{} \ill{} \uncertain{injectif}, \ill{}
\[
  \Omega^{1}_{k \cap K^{p}} \otimes_{k \cap K^{p}} K \to
  \Omega^{1}_{k} \otimes_{k} K \to \Omega^{1}_{K}
\]
\note{le second $K$ de cette suite est biffé sur le feuillet}
\ill{} \ill{} \uncertain{l'existence} \uncertain{dessus} \ill{}
\ill{} il \ill{} \ill{} \uncertain{avoir} \ill{} $N_{K/k}$ \ill{}
\[
  \Omega_{k} \otimes_{k} K \longrightarrow \Omega_{K}
\]
\ill{} \ill{} \ill{} \uncertain{$d x \otimes 1$} \quad
($x \in k \cap K^{p}$).

N.B. \struck{\ill{}} \ill{} Voici un contre-exemple. \ill{} $N_{K/k}$
\ill{} \ill{} \ill{} Prenons \ill{} \ill{} \ill{} premier \ill{} \ill{}
\ill{} \ill{} les \ill{} déterminations \ill{}

\struck{$K^{p} = \mathbf{F}(x^{p}, y^{p}, (a_i))$} \qquad
\struck{$K = \mathbf{F}(x, y, (a_i))_{i \geq 1}$}

N.B. $a$ et $(a_i)_{1 \leq i \leq p-1}$ \ill{} \ill{} alg.\ \ill{} \ill{}
$\mathbf{F}$

\[
  \begin{cases}
    a_0 = y - \sum_{1}^{p-1} a_i x^{i} \\
    a = x^{p}
  \end{cases}
  \qquad \text{d'où} \qquad y = \sum_{0}^{p-1} a_i x^{i}
\]
\marginal{$(1 \leq i \leq p-1)$}
\[
  K^{p} = \mathbf{F}(a, a_0, \dots, a_p)
  = \struck{\mathbf{F}(\alpha, x^{p}, y - \alpha x)}\;
    \mathbf{F}\Big(x^{p}, a_1, \dots, a_{p-1},\, y - \sum_{1}^{p-1} a_i x^{i}\Big)
\]
\note{la première expression du membre de droite est biffée et remplacée par la
seconde, écrite au-dessus}
\[
  \struck{K = \mathbf{F}\big(\alpha^{1/p}, x, (y - \alpha x)^{1/p}\big)}
  \qquad \text{i.e.} \qquad
  K = \mathbf{F}\big(x, a_1^{1/p}, \dots, a_{p-1}^{1/p}, \ill{}\big)
\]
$\mathbf{F}(x,y)$ \quad [N.B. $x, y \in K$, \ill{}
$y = \alpha x + \beta \in k$, \ill{} $k \subset K$ \ill{}
\note{le crochet ouvert de ce N.B. n'est pas refermé : la fin de la ligne se
perd dans le bas du feuillet}

\page{28}
\struck{\ill{} \ill{} $x, y \notin$ \ill{}} \uncertain{Voir} \ill{}
\ill{}, \ill{} $x \in K$, $x \notin K^{p}$, \ill{}
\marginal{(cf. le N.B. \uncertain{précédent})}
$y \notin K^{p}$, $y \in K^{p}(x)$, donc \ill{} les $a_i$ \ill{} \ill{}
\ill{} \ill{} \emph{\ill{}} $\in K^{p}$. Comme $x \notin K^{p}$, \ill{}
\ill{} $y = \sum_{0}^{p-1} a_i x^{i}$.
\struck{$y = \alpha x + \beta$. Évident}

\ill{} $x\,y \in K^{p} \cap k$, \ill{} $x \notin K^{p} \cap k$, mais bien
\ill{} \ill{} \ill{} $\cap\, k$, si donc $(x,y)$ n'était \struck{\ill{}}
\ill{} $K^{p} \cap k$, \ill{} \ill{} \ill{}
$y = \sum_{0}^{p-1} a_i x^{i}$, les $a_i \in K^{p} \cap k$, donc
\ill{} \ill{} les $a'_i = a_i$, \ill{} \ill{} donc
$a_i \in k = \mathbf{F}(x,y)$, absurde.

N.B. Condition \uncertain{isomorphes} \ill{}
\[
  \begin{cases}
    (x,y) \;\text{p-libre sur}\; \ill{} \\
    \ill{}\; \lambda \notin \mathbf{F}(x,y)
  \end{cases}
\]
\struck{Je \uncertain{voudrais} \ill{}} exemple
\struck{$\lambda = \mathbf{F}(a_1, \dots, a_{p-1})$, et \ill{}}
\struck{(par $x, y$ \ill{} $\mathbf{F}(x,y)$ \ill{} \ill{})}
\struck{\ill{} par \ill{} \ill{} $\lambda$ \ill{} $a_i$ \ill{}}
\ill{} $\mathbf{F}(x,y,a)$, \ill{}
\[
  \lambda = \mathbf{F}(a_1, \dots, a_{p-1})
\]
Prenons $x, y, a_1, \dots, a_{p-1}$ alg.\ indép.\ sur $\mathbf{F}_0$
\add{et prendre} $a_i = 0$ pour $2 \leq i \leq p-1$.

\ill{} \ill{} \ill{} \ill{} \ill{} Bien entendu, $\mathbf{F}$ peut
\ill{} \ill{} peut aussi \ill{} être un corps fini \ill{} \ill{} \ill{}
\ill{} $x, y$ alg.\ indép.\ sur $\mathbf{F}$, \ill{} et $a_1$
\struck{\ill{}} sur $\mathbf{F}$, $a_1 \notin \mathbf{F}$, \ill{} $K$
\ill{} obtenu \ill{} \ill{} \ill{} \ill{} ($y - \alpha x$) \ill{}
$k(b)$, où $b$ \ill{} \ill{} $\mathbf{F}$ \quad $[\,b = a^{1/p}\,]$.

\end{document}
